% =============================================================================
% Chapter 6 — Evaluation of Intent-Conditioned Single-Path Selection
%
% Reframed as a MOTIVATING STUDY for Part II: an intent-conditioned Q-network
% (FiLM) selects near-optimal single paths at low probing cost, but hits an
% architectural ceiling under congestion/volatility — motivating the joint
% multipath + bitrate co-optimization of Part II.
%
% Requires (in the thesis preamble):
%   \usepackage{graphicx,booktabs,subcaption,amsmath,xcolor,siunitx}
% Figures + table are produced by:  evaluation/run_chapter6.py <run_dir>
% then copy the produced  chapter6_<ts>/  contents into docs/chapter6/
% (table_6_1.tex and figures/*.png), matching \chsixdir below.
% =============================================================================

% ---- headline numbers (populated from the chapter-6 evaluation run) ----------
% Source: run_20260722_180329/chapter6_20260722_200655
\newcommand{\chsixdir}{chapter6}          % dir holding table_6_1.tex + figures/
\newcommand{\filmDivergence}{0.128}       % Conditional-FiLM adaptivity
\newcommand{\concatDivergence}{0.050}     % value-only Conditional-Concat adaptivity
\newcommand{\flatDivergence}{0.011}       % Flat DQN adaptivity (noise floor)
\newcommand{\probeReductionPct}{97.5\%}   % FiLM vs. widest-path exhaustive probing
\newcommand{\filmProbeMs}{\SI{135}{\milli\second}}
\newcommand{\widestProbeMs}{\SI{5388}{\milli\second}}
\newcommand{\filmGoodputHi}{9067\,Mbps}   % lowest-congestion bin
\newcommand{\filmGoodputLo}{6844\,Mbps}   % highest-congestion bin
\newcommand{\latencyIntentMs}{\SI{12.8}{\milli\second}}  % chosen latency under Low-Latency
\newcommand{\otherIntentMs}{\SI{16}{\milli\second}}      % chosen latency under other intents
% -----------------------------------------------------------------------------

\chapter{Evaluation: Intent-Conditioned Single-Path Selection and its Ceiling}
\label{ch:evaluation}

This chapter closes Part~I and sets up Part~II. Its question is deliberately
narrow: \emph{how well can a single agent, told an operator intent, choose one
path for a flow?} We show that a Q-network which conditions on the intent through
\emph{feature-wise linear modulation} (FiLM) learns to re-rank paths as the
intent changes, matching or beating hand-tuned heuristics while probing only the
path it actually selects. We then show the limit of the single-path formulation:
as the network becomes congested and volatile, even a perfectly-conditioned
single-path choice degrades, because no single path can simultaneously satisfy
the throughput a flow demands. That observation is the bridge to Part~II, where
the agent is allowed to \emph{split} traffic across paths and co-optimize the
application bitrate.

All results use the SCION AS-level simulator of Chapter~\ref{ch:simulator}. The
agents are trained on the first 14 simulated days and evaluated on the held-out
last 14 days over up to 32 routable source--destination pairs, each offering 30
candidate paths. Reward is the composite goodput/trust score of
Section~\ref{sec:reward}, $r = 2\,(w_1 G + w_2 T) - 1 - w_{\text{probe}}\,c$,
where $G$ is goodput normalized to a per-topology cap, $T$ a link-trust term
whose internal shape is set by loss/delay weights $(w_3,w_4)$, and $c$ a
normalized probing penalty. Four \emph{intents} instantiate this reward with
different weight vectors and define genuinely distinct selection objectives:
\textbf{Throughput} ($w_1\!\to\!1$), \textbf{Low-Latency} (trust dominated by the
delay term, $w_4\!\to\!1$), \textbf{Low-Loss} (trust dominated by the loss term,
$w_3\!\to\!1$), and \textbf{Balanced}. Probe-cost preference ($w_{\text{probe}}$)
is treated separately in Section~\ref{sec:ch6-probing}: for a selective agent
that probes a single path, probing cost is independent of \emph{which} path is
chosen, so it is not an intent that can change path ranking.

\section{Agents and the conditioning question}
\label{sec:ch6-agents}

We compare a non-conditioned reference against two ways of making a
path-scoring network intent-aware. All three share the same dueling path-scoring
backbone; they differ only in \emph{where} the intent vector $w$ enters:
\begin{itemize}
  \item \textbf{Flat DQN} — a masked value network over a global state that does
  not receive $w$ at all. The non-conditioned reference: one policy, one path
  per state.
  \item \textbf{Conditional-Concat (value-only)} — $w$ is concatenated into the
  \emph{value} stream only; the per-path advantage stream never sees it. This is
  the naive way to add an intent input.
  \item \textbf{Conditional-FiLM} — $w$ is mapped to per-path scale/shift
  parameters $(\gamma,\beta)$ that \emph{modulate} each path's features in the
  advantage stream, while the value stream sees only topology/time context.
\end{itemize}

The distinction is not cosmetic; it follows from the dueling decomposition
$Q(s,\text{path}_i) = V(s) + \big(A(s,\text{path}_i) - \bar{A}(s)\big)$. The
selected path is $\arg\max_i A(s,\text{path}_i)$, which does not depend on
$V(s)$. Hence \emph{intent can only change the chosen path if it reaches the
advantage stream.} Injecting $w$ into the value stream alone (value-only concat)
shifts $V(s)$ by a constant that is identical across paths and cancels in the
$\arg\max$: such a network is \emph{structurally unable to re-rank paths} as the
intent changes.\footnote{We verify this property directly at the network level:
holding the path features fixed and varying only $w$ leaves the per-path
advantage---and therefore the selected path---unchanged, while $V(s)$ shifts
uniformly.} FiLM, by construction, lets $w$ act on the per-path advantages, so a
change of intent \emph{can} change which path wins.
Section~\ref{sec:ch6-ablation} makes this concrete.

\section{Ablation: does the network actually adapt to intent?}
\label{sec:ch6-ablation}

Table~\ref{tab:ch6-ablation} reports the mean intent-weighted reward of each
agent under each intent, together with an \emph{adaptivity} score. Adaptivity is
the \emph{behavioral divergence}: for every (pair, hour) context we record the
path each agent selects under all four intents and measure the fraction of
contexts in which the choice actually changes, normalized to $[0,1]$ (0 = the
agent ignores the intent when choosing a path; 1 = it re-ranks maximally).

\begin{table}[t]
  \centering
  \caption{Mean intent-weighted reward by agent and intent on the held-out
  window, with the behavioral-divergence (adaptivity) score. Conditioning lifts
  reward far above the non-conditioned Flat DQN under every intent. Flat DQN and
  the value-only concat leave their path ranking essentially unchanged across
  intents (adaptivity $\approx 0$), whereas FiLM---which routes the intent into
  the per-path advantage stream---re-ranks paths as the intent vector changes.}
  \label{tab:ch6-ablation}
  % Auto-generated by src/pipeline/chapter6_eval.py -- do not edit by hand.
\begin{tabular}{lccccc}
\toprule
Method & Throughput & Low-Latency & Low-Loss & Balanced & Adaptivity \\
\midrule
Flat DQN & 0.333 & 0.701 & 0.868 & 0.561 & 0.011 \\
Conditional-Concat & 0.979 & 0.708 & 0.937 & 0.894 & 0.050 \\
Conditional-FiLM & 0.980 & 0.732 & 0.939 & 0.894 & 0.128 \\
\bottomrule
\end{tabular}

\end{table}

Two results stand out. First, \emph{conditioning works}: both the FiLM and the
concat agents reach a mean reward around $0.9$--$0.98$ on the Throughput,
Low-Loss and Balanced intents, against $0.33$--$0.87$ for the non-conditioned
Flat DQN---the intent input is being used. Second, and decisively, \emph{whether
the network re-ranks paths depends entirely on where the intent enters.} The
Flat DQN (adaptivity $\flatDivergence$) and the value-only concat (adaptivity
$\concatDivergence$) sit at the non-adaptive floor: telling the value-only concat
a different intent almost never changes the path it picks, exactly as the
dueling decomposition predicts.\footnote{The small residual $\concatDivergence$
for the value-only concat is not the conditioning mechanism at work: one path
feature (the per-path trust score) is itself a function of $(w_3,w_4)$, so a
sliver of intent leaks into the advantage stream through the \emph{features}
rather than the architecture. The value stream contributes nothing to
re-ranking, as the network-level test confirms.} In contrast, FiLM
(adaptivity $\filmDivergence$)---whose intent reaches the per-path advantage
stream---re-ranks paths several times more often, at no cost in reward: it
matches the value-only concat's reward while conditioning through a compact
per-path modulation. Conditioning turns an intent \emph{input} into
intent-\emph{driven behavior} only when it modulates the per-path advantages.

\section{Intent alignment and behavioral shift}
\label{sec:ch6-alignment}

Having established that FiLM adapts, we characterize \emph{how} it adapts.
Figure~\ref{fig:ch6-heatmap} evaluates the FiLM agent under every combination of
the intent it is \emph{told} (rows) and the objective its chosen path is
\emph{scored} under (columns). Because each objective sits at its own intrinsic
reward level, each cell is coloured by its \emph{deviation from that column's
mean} on a single shared, diverging scale: red where a conditioning intent beats
the average intent for that objective, blue where it does worse. Colour intensity
is thus proportional to the true reward gap, and a well-aligned agent puts the
red on the (outlined) diagonal.

\begin{figure}[t]
  \centering
  \includegraphics[width=0.66\textwidth]{\chsixdir/figures/fig_6_1_heatmap.png}
  \caption{Reward matrix $R(\text{intent told}, \text{objective scored})$ for
  Conditional-FiLM. Cells are annotated with the raw mean reward and coloured by
  their deviation from the column (objective) mean; the outlined diagonal marks
  the matched intent. The Throughput and Low-Latency objectives are maximized by
  conditioning on their own intent (diagonal is the reddest cell in its column),
  while the Low-Loss and Balanced columns are near-flat---in this topology the
  candidate paths barely differ in loss, so conditioning has little to re-rank.
  Where path choice matters, telling the agent an intent moves its choice toward
  that objective.}
  \label{fig:ch6-heatmap}
\end{figure}

The mechanism behind the diagonal is a genuine shift in \emph{which} paths the
agent picks. Figure~\ref{fig:ch6-boxplots} shows the distribution of the chosen
path's latency, goodput, and trust as the conditioning intent varies. Under the
Low-Latency intent the agent selects markedly lower-latency paths---median
$\latencyIntentMs$ against $\approx\otherIntentMs$ under the other intents---and
accepts lower goodput in exchange, the correct trade-off. Under the Throughput
intent it selects the highest-goodput paths, and under Low-Latency it also
attains the highest chosen-path trust. All of this comes from a single set of
weights, selected at inference time by the intent vector alone.

\begin{figure}[t]
  \centering
  \includegraphics[width=\textwidth]{\chsixdir/figures/fig_6_1_boxplots.png}
  \caption{Distribution of the chosen path's latency, goodput, and trust across
  the four conditioning intents (Conditional-FiLM). The agent moves its
  selections toward whichever metric the active intent rewards---most visibly,
  the Low-Latency intent shifts the whole latency distribution down while trading
  away goodput.}
  \label{fig:ch6-boxplots}
\end{figure}

\section{Probing overhead: quality at a fraction of the cost}
\label{sec:ch6-probing}

A path selector is only useful if the information it needs is cheap to obtain.
The heuristics (Shortest-Path, Widest-Path, Lowest-Latency, Random) are
\emph{exhaustive}: to choose, they must probe \emph{every} candidate path---a
latency probe ($\SI{10}{\milli\second}\,+\,\SI{0.5}{\milli\second}$/hop) or, for
bandwidth-aware heuristics, a full probe
($\SI{100}{\milli\second}\,+\,\SI{20}{\milli\second}$/hop) per path. The FiLM
agent is \emph{selective}: it probes only the single path it decides to use.

Figure~\ref{fig:ch6-probe} plots selection quality against probing cost per
selection. Conditional-FiLM sits in the desirable upper-left corner: it achieves
essentially the same goodput as the best exhaustive heuristic (Widest-Path,
8242\,Mbps) at \filmProbeMs{} of probing per selection versus
\widestProbeMs{}---a reduction of \probeReductionPct{}---and in composite reward
it slightly \emph{beats} Widest-Path because it is not charged for probing the 29
paths it does not use. The latency-probe heuristics reach only
5800--6100\,Mbps despite probing all 30 paths.

\begin{figure}[t]
  \centering
  \includegraphics[width=0.64\textwidth]{\chsixdir/figures/fig_6_2_quality_vs_probe.png}
  \caption{Selection quality (achieved goodput) versus probing cost per
  selection. Conditional-FiLM achieves the best heuristic's quality while probing
  only its chosen path; the bandwidth-aware Widest-Path pays roughly $40\times$
  the probing cost to reach comparable quality.}
  \label{fig:ch6-probe}
\end{figure}

\section{The single-path ceiling}
\label{sec:ch6-ceiling}

The central limitation of Part~I appears when we stratify performance by realized
congestion. Figure~\ref{fig:ch6-ceiling} plots achieved goodput against the mean
utilization of the flow's candidate paths at the moment of selection. At low
congestion, path choice is nearly a free lunch: FiLM and the bandwidth-aware
Widest-Path both deliver around \filmGoodputHi{}. As utilization rises, \emph{every}
method's goodput falls: FiLM declines to \filmGoodputLo{} at the highest
congestion, and the delay- and hop-based heuristics collapse faster still.
Conditional-FiLM tracks the best heuristic at every congestion level---its curve
sits on top of Widest-Path throughout, at a fraction of the probing cost---but it
cannot escape the decline. When the good paths are saturated, choosing the best
\emph{single} path still leaves the flow starved; no amount of smarter
\emph{single}-path selection closes that gap.

\begin{figure}[t]
  \centering
  \includegraphics[width=0.68\textwidth]{\chsixdir/figures/fig_6_3_ceiling.png}
  \caption{Application performance (achieved goodput) versus realized network
  congestion. As congestion rises, all single-path selectors degrade; the best
  single path is no longer good enough. This ceiling is intrinsic to single-path
  selection, not to the policy---Conditional-FiLM already tracks the best
  heuristic everywhere.}
  \label{fig:ch6-ceiling}
\end{figure}

\section{Summary and bridge to Part~II}
\label{sec:ch6-summary}

Part~I establishes two things. First, \emph{intent conditioning works, but only
when the intent reaches the per-path advantage stream}: FiLM modulation turns the
intent vector into genuine, measurable re-ranking of path selections (adaptivity
$\filmDivergence$ versus $\flatDivergence$ for the flat baseline and
$\concatDivergence$ for a value-only concat that is architecturally blind to
intent in its ranking), producing intent-appropriate behavior---lower-latency
paths on demand, highest-goodput paths on demand---at a probing cost far below
exhaustive heuristics. Second, and decisively for what follows, \emph{single-path
selection has a ceiling}: under congestion the achievable goodput falls from
\filmGoodputHi{} to \filmGoodputLo{} even for the best-conditioned policy,
because the bottleneck is the commitment to one path, not the quality of the
choice.

Part~II removes exactly that commitment. Instead of selecting one path per flow,
the agent learns to \emph{split} traffic across multiple paths and to
\emph{co-optimize the application bitrate} against the resulting aggregate
capacity. The intent-conditioning machinery validated here---FiLM modulation of
per-path features---carries over intact; what changes is the action space, and
with it the ceiling.
